\subsection{The exact elliptic cover and the remaining relative automorphism group}

\begin{statement}{D.34}{Proved}{A marked elliptic-cover and centralizer description of the relative pro-two source}
This addresses Part D(e) and the CM follow-up to LM2, manuscript
\S51.16, folio 734 verso, scan 174, and \S51.15, folios 725--727,
scans 156--160. Retain $B=\langle s,r\mid s^2=r^3=1\rangle$,
$t=sr$, $Q_0$, $K$ and $\mathcal P_2$ from D.30.
The index-18 subgroup $K$ is the fundamental group of
\[
 \mathcal E_0^\circ:\quad w^3=z(z-1),\qquad
 z\notin\{0,1,\infty\}.                              \tag{D.34.1}
\]
Its smooth compactification has genus one. Its three removed points
are denoted $P_0,P_1,P_\infty$. The regular $Q_0$ dessin has
permutation cycles $2^9,3^6,6^3$, and $Q_0\simeq S_3\times C_3$ acts
by the following explicit deck maps. For
$\zeta=\exp(2\pi i/3)$ set
\[
 \delta(z,w)=(z,\zeta w),\quad
 s(z,w)=(1-z,w),\quad
 r_0(z,w)=\left(\frac1{1-z},\frac{-w}{1-z}\right),
 \qquad r=\delta r_0.                                 \tag{D.34.2}
\]
Here $s^2=r_0^3=(sr_0)^2=1$, $\delta$ is central,
$t^2=\delta^2$, and the standard positive generators
$x=t^{-2}$ and $y=sxs$ of $\Gamma(2)$ both act as $\delta$.

In the fixed rank-four basis $a_1,a_2,a_3,a_4$ of D.30, the three
cusp generators, corresponding respectively to $z=0,1,\infty$, are
\[
 c_0=a_1a_4^{-1}a_2=t^6,\qquad
 c_1=a_1^{-1}a_3a_2^{-1},\qquad
 c_2=a_3^{-1}a_4.                                     \tag{D.34.3}
\]
The negative choice of orientation implicit in $t^6=x^{-3}$ is the
same at all three cusps. Their full $Q_0$ conjugacy maps are specified
in the proof and in the 54-entry word certificate.

Put $K_2=K^{(2)}$ and let
$\Lambda\simeq Q_0\subset\operatorname{Out}(K_2)$ be the outer
conjugation action. Restriction is an isomorphism of profinite groups
\[
 \operatorname{Aut}_{Q_0}(\mathcal P_2)
 \xrightarrow{\ \sim\ }
 \{a\in\operatorname{Aut}(K_2):\bar a\in
                C_{\operatorname{Out}(K_2)}(\Lambda)\}.          \tag{D.34.4}
\]
The subscript on the left means that the induced map on the displayed
quotient $Q_0$ is the identity. Under this isomorphism, the exact
normalized boundary condition is
\[
 a(c_0)=c_0^{\lambda_2}
 \quad\Longleftrightarrow\quad
 \alpha(t)=t^{(\lambda_2,1)},\qquad
 \lambda_2\in\mathbb Z_2^\times,                     \tag{D.34.5}
\]
where $\overline{\langle t\rangle}\simeq
\mathbb Z_2\times C_3$, and the two entries are its two-adic and
order-three exponents. Both sides then send every $c_i$ to a
$K_2$-conjugate of $c_i^{\lambda_2}$.

Filling the three cusps gives the exact map
\[
 K_2\longrightarrow T_2\mathcal E_0,
 \qquad
 0\longrightarrow(\mathbb Z_2^2)_{\rm cusp}
 \longrightarrow H_1(K_2,\mathbb Z_2)
 \longrightarrow T_2\mathcal E_0\longrightarrow0.     \tag{D.34.6}
\]
The sequence has a canonical $Q_0$-equivariant splitting over
$\mathbb Z_2$, using the central order-three idempotent. On its
rank-two elliptic factor, $t^2$ has matrix
\[
 Z=\begin{pmatrix}-1&1\\-1&0\end{pmatrix},\qquad
 Z^2+Z+I=0.                                           \tag{D.34.7}
\]
Every automorphism in (D.34.5) therefore acts by
$\lambda_2 I$ on the cusp factor and by a unit
$a+bZ\in\mathbb Z_2[\zeta]^\times$ on the elliptic factor, with
\[
                    a^2-ab+b^2=\lambda_2.             \tag{D.34.8}
\]
The retained $T_3$ component of D.33 has the following exact reduction
on the latter factor:
\[
 M_{16}=\begin{pmatrix}3&8\\8&11\end{pmatrix}
        =11I+8Z\pmod {16},\qquad\det M_{16}=1\pmod {16}.
                                                               \tag{D.34.9}
\]
All the stable corrections of D.33 retain this matrix.

For the arithmetic comparison, let $c=\sqrt[3]{4}>0$,
$E_3:Y^2=X^3+1$, and $P=(0,1)\in E_3[3]$. The exact isomorphism
with origin $P_0$ on $\mathcal E_0$ is
\[
 \phi=T_P\circ\phi_0,\qquad
 \phi_0(z,w)=(cw,2z-1),\qquad
 \phi:\mathcal E_0\longrightarrow E_3.              \tag{D.34.10}
\]
It is defined over $\mathbb Q(c)$. If $\Phi=T_2\phi$ and
$\kappa_{4,3}(g)$ is determined by $g(c)/c=\zeta^{\kappa_{4,3}(g)}$,
then the actual Galois representations satisfy
\[
 \Phi\rho_{\mathcal E_0}(g)\Phi^{-1}
       =[\zeta]^{-\kappa_{4,3}(g)}\rho_{E_3}(g),
 \qquad\kappa_{4,3}=2\kappa_{2,3}.                   \tag{D.34.11}
\]
Thus this comparison retains the cubic Kummer twist. Equation
(D.34.11) concerns actual Galois action; (D.34.4)--(D.34.9) give the
specified maps for the larger relative automorphism group.
\end{statement}
\begin{proof}
The program computes the regular right-multiplication permutations of
$s,r,t$ on the 18 explicitly listed elements of $Q_0$. Their respective
cycle lengths are nine twos, six threes and three sixes. The compact
dessin therefore has Euler characteristic $9+6-18+3=0$, and genus one.
Removing its three faces leaves the graph of Euler characteristic
$-3$ whose four generators were retained in D.30.

There is also an algebraic identification retaining the actual
monodromy character. The homomorphism
$Q_0\to S_3\times C_3$ sends a word to its action on the three
$\Gamma(2)$ cusps and to its exponent sum in $r$ modulo three.
All 18 pairs are distinct, and multiplication is checked on all
$18^2$ pairs. Thus $K$ is the kernel on $\Gamma(2)$ of the character
sending $x=t^{-2}$ and $y=sxs$ to $1\in\mathbb Z/3$.
For the map (D.34.1) to the $z$ line, a positive loop about either
zero or one multiplies $w$ by $\zeta$. The third positive loop does
the same, since $z(z-1)$ has pole order two at infinity. This is
exactly the indicated character, proving the marked cover comparison.
The full induced pro-two topology is the one proved in D.30.

The maps $s,r_0$ in (D.34.2) preserve $w^3=z(z-1)$ and satisfy the
displayed $S_3$ relations by direct substitution. They commute with
$\delta$. Their 18 products are distinct, since the six base maps
are distinct and each has three different scalar lifts. The invariant
map to the orbifold line is
\[
 j_0(z)=\frac{4(z^2-z+1)^3}{27z^2(z-1)^2}.
                                                               \tag{D.34.12}
\]
It is invariant under both base generators. Consequently these maps
realize the exact $Q_0$ cover of the orbifold, with the stated $s,r$
and cyclic character.

For completeness every cusp map is retained at the word level. Choose
the recorded short $B$ representatives $g_i$ of the three left cosets
of $\langle t\rangle$ in $Q_0$, with $g_0=1$. Then
$c_i=g_it^6g_i^{-1}$ gives (D.34.3). For a representative $h$ of
$q\in Q_0$, write $qg_i=g_jt^a$ in $Q_0$, $0\leq a<6$, and put
\[
 k_{q,i}=hg_it^{-a}g_j^{-1}\in K.
 \quad\text{Then}\quad hc_ih^{-1}=k_{q,i}c_jk_{q,i}^{-1}.
                                                               \tag{D.34.13}
\]
The certificate records all 54 resulting $K$ words. Each equality
is independently reduced in $C_2*C_3$. It also records all section
cocycles $h_qh_{q'}h_{qq'}^{-1}\in K$ and verifies their conjugation
identity on all four $K$ generators, giving 1296 exact checks. This
specifies the outer $Q_0$ action, including its extension data.

The 18 matrices of that action on $H_1(K,\mathbb F_2)$ are distinct.
Thus an element of $\mathcal P_2$ centralizing $K_2$ has trivial
image in $Q_0$, and then lies in the center of the nonabelian free
pro-two group $K_2$, which is trivial. Conjugation consequently embeds
$\mathcal P_2$ into $\operatorname{Aut}(K_2)$. Its image is exactly
the inverse image of $\Lambda$ under
$\pi:\operatorname{Aut}(K_2)\to\operatorname{Out}(K_2)$: it contains
every inner automorphism and one representative of every member of
$\Lambda$.

Restriction now gives (D.34.4) in both directions. An automorphism of
$\mathcal P_2$ over $Q_0$ restricts to an automorphism whose outer class
centralizes $\Lambda$. Conversely, conjugation by any such
$a\in\operatorname{Aut}(K_2)$ preserves $\pi^{-1}(\Lambda)$ and
induces the identity on $\Lambda$. Transporting this conjugation
through the preceding embedding gives the required automorphism of
$\mathcal P_2$. The restriction is exactly $a$, because
$a\operatorname{Int}(k)a^{-1}=\operatorname{Int}(a(k))$.
Faithfulness makes these operations inverse. Compactness makes the
displayed continuous bijections isomorphisms of profinite groups.

We verify the boundary condition without assuming a root uniqueness
statement for the whole group. The stabilizer in $Q_0$ of
$[c_0]=(1,1,0,1)$ in $H_1(K,\mathbb F_2)$ is precisely
$\langle t\rangle$, of order six. The element $c_0$ is primitive in
$K_2$: (D.34.3) permits replacing $a_4$ by $c_0$, with inverse
$a_4=a_2c_0^{-1}a_1$. The free-product centralizer theorem
\cite[Theorem B$'$]{A_HerfortRibes1985} therefore gives
$C_{K_2}(c_0)=\overline{\langle c_0\rangle}$.
For any element centralizing $c_0$, its $Q_0$ image is a power of $t$;
remove that power and use this last equality. We obtain
\[
 C_{\mathcal P_2}(c_0)=\overline{\langle t\rangle}
                      \simeq\mathbb Z_2\times C_3.    \tag{D.34.14}
\]
The procyclic description follows because its kernel over its order-six
image is $\langle t^6\rangle\simeq\mathbb Z_2$. If
$a(c_0)=c_0^{\lambda_2}$, then $\alpha(t)$ centralizes $c_0$, since
it centralizes $c_0^{\lambda_2}$ and this element generates the same
closed subgroup. Inside (D.34.14), its sixth power fixes its two-adic
exponent as $\lambda_2$, while the unchanged $Q_0$ image fixes its
order-three exponent as $1$. This proves (D.34.5), including the
converse. Equation (D.34.13), and the fact that $\alpha(h)h^{-1}\in K_2$,
then prove the assertion at every cusp.

Filling the punctures is an actual morphism of curves. On topological
fundamental groups its kernel is the normal closure of the $c_i$.
Its compact genus-one quotient is $\mathbb Z^2$, so pro-two completion
gives (D.34.6). The identification with $T_2\mathcal E_0$ uses the
cofinal system of covers $[2^n]$ of this pointed elliptic curve. In the
displayed $K$ basis the integral filling matrix and a section are
\[
 P=\begin{pmatrix}1&0&1&1\\0&1&1&1\end{pmatrix},
 \qquad
 J=\begin{pmatrix}1&0\\0&1\\0&0\\0&0\end{pmatrix}.
                                                               \tag{D.34.15}
\]
The cusp lattice is generated by $(1,1,0,-1)$ and $(-1,-1,1,0)$.
These, followed by the two columns of $J$, form an integral basis.
Thus $P$ is the precise quotient map, rather than a rational
identification of lattices.

Let $z_0=t^2$. The matrix $P M_{z_0}J$ is (D.34.7). On the rank-four
lattice the idempotent $(1+z_0+z_0^2)/3$ has image exactly the cusp
lattice; its complementary image maps isomorphically by $P$ to the
elliptic quotient. This proves the canonical splitting. The latter
matrix has irreducible polynomial $X^2+X+1$ modulo two. A vector and
its $Z$ translate are therefore a basis over $\mathbb Z_2$, so its
centralizer ring is exactly $\mathbb Z_2[Z]$. Its determinant is
$a^2-ab+b^2$.

The determinant equals the common cusp multiplier. One can see this
directly from the fundamental class of the punctured torus relative
to its boundary circles. Its boundary is the sum of the three oriented
boundary classes, so a common power $\lambda_2$ acts by $\lambda_2$
on that relative fundamental class. Naturality of the intersection
pairing on the compact rank-two quotient then gives
$M^{\mathsf T}\left(\begin{smallmatrix}0&1\\-1&0\end{smallmatrix}\right)M
=\lambda_2\left(\begin{smallmatrix}0&1\\-1&0\end{smallmatrix}\right)$.
For a two-by-two matrix this is exactly (D.34.8). The same relative
pair and its orientation are covered by
\cite[Theorem 3.3 and \S1.3]{dWilkes2020}; here the integral filling
map has been given explicitly.

For (D.34.9), the program reconstructs the exact D.33 conjugators.
It calculates their action on all four $K$ generators in abelianization
modulo sixteen, obtaining
\[
 M_K=\begin{pmatrix}
 7&12&2&2\\12&15&10&10\\14&6&5&4\\14&6&4&5
 \end{pmatrix}\pmod {16}.
                                                               \tag{D.34.16}
\]
This matrix commutes with every $Q_0$ matrix modulo sixteen and fixes
both cusp vectors. Multiplication $P M_KJ$ gives (D.34.9).
The quotient $T_3$ retains $H_1(K,\mathbb Z/16)$ because the
abelian image of $D_9(K)$ is $16\mathbb Z_2^4$. D.33 retains its
entire $T_3$ action, so it retains this Tate component as well.

Finally (D.34.10) follows by substituting
$(2z-1)^2=4z(z-1)+1=(cw)^3+1$; its inverse before translation is
$z=(Y+1)/2$, $w=X/c$. The point $P$ has order three: its horizontal
tangent gives $2P=-P$. Under $\phi_0$, $\delta$ is $[\zeta]$,
$s$ is $[-1]$, and $r_0$ is translation by $-P$. After translating
the origin by $P$, the deck maps become
\[
       s=T_{-P}[-1],\qquad r=T_{-P}[\zeta],
       \qquad t=[-\zeta].                             \tag{D.34.17}
\]
These are the claimed translations of order three combined with
$\mu_6$. Substitution in (D.34.12) gives the Belyi map
$(X^3+4)^3/(27X^6)$, invariant under these same deck transformations.
For $g\in G_{\mathbb Q}$, the exact equality of morphisms
$g(\phi)=[\zeta]^{\kappa_{4,3}(g)}\phi$ holds: $P$ is rational and
is fixed by $[\zeta]$. Functoriality on Tate modules proves
(D.34.11). The equality $\kappa_{4,3}=2\kappa_{2,3}$ follows from
$c=(\sqrt[3]{2})^2$. This proves the arithmetic comparison with its
specified field of definition and twisting character.
\end{proof}
